\section{Marco teórico}

\subsection{Termopares}
Un termopar (o cupla térmica) es un sensor de temperatura industrial que se compone de dos hilos de metales diferentes unidos en un extremo. Cuando existe una diferencia de temperatura entre la unión de los metales y los extremos libres, se genera una fuerza electromotriz que es proporcional a esa variación.
Según el material del que están compuestos, cambia el rango de temperaturas que pueden medir. 

\subsection{Termistores}

Los termistores son dispositivos semiconductores cuya resistencia eléctrica varía en función de la temperatura. Según el sentido de esta variación, pueden clasificarse en dos grupos: NTC (\textit{Negative Temperature Coefficient}) y PTC (\textit{Positive Temperature Coefficient}).

En uno NTC, la resistencia disminuye conforme aumenta la temperatura. Por el contrario, en un termistor PTC, la resistencia aumenta con la temperatura. En ambos casos, la relación inversa entre la resistencia y la temperatura permite utilizarlos como sensores. Sin embargo, su respuesta frente a la temperatura es no lineal. Por ello, para obtener la temperatura a partir de una resistencia medida es necesario utilizar un modelo matemático que describa esta relación.


\subsection{Ecuación de Steinhart-Hart}

La ecuación de Steinhart-Hart permite relacionar la resistencia eléctrica de un termistor con su temperatura absoluta. Su expresión general es:

\begin{equation}
\frac{1}{T}
=
A+B\ln(R)+C\left[\ln(R)\right]^3,
\end{equation}

donde $T$ es la temperatura absoluta expresada en kelvin, $R$ es la resistencia eléctrica del termistor expresada en ohmios y $A$, $B$ y $C$ son los coeficientes de Steinhart-Hart.

Despejando la temperatura de la expresión anterior:

\begin{equation}
T
=
\frac{1}
{A+B\ln(R)+C\left[\ln(R)\right]^3}.
\end{equation}

Los coeficientes $A$, $B$ y $C$ dependen de las características del termistor y pueden obtenerse a partir de tres pares de valores de resistencia y temperatura. Para ello, se utilizan tres mediciones realizadas a diferentes temperaturas y se resuelve el sistema de ecuaciones correspondiente.

